\section*{Capítulo 13}

\solucion{13.1} Ejemplo: una aerolínea compra hoy un forward de combustible para un consumo de dentro de seis meses. Al cierre del trimestre el combustible ha bajado y el forward tiene valor negativo de 5 millones, que se registra en resultados. El ahorro en el combustible que se comprará más barato no se registra todavía. La utilidad del trimestre cae 5 millones aunque, económicamente, la empresa sigue pagando el mismo precio fijo. En el trimestre siguiente, cuando se consuma el combustible barato, la utilidad sube. El derivado redujo el riesgo económico pero trasladó volatilidad entre trimestres.

\solucion{13.2} El saldo en dólares del préstamo baja: $3{,}765.6/3.7817-3{,}765.6/3.4379=995.7-1{,}095.3=-99.6$, una ganancia cambiaria de 99.6 millones en resultados. El CCS pasa a ser un pasivo de 101.9 millones; su cambio de valor se registra en ORI ($-101.9$). Se reclasifican de ORI a resultados $-99.6$ millones, que compensan la ganancia del préstamo. El resultado del período queda neutral y en ORI quedan unos $-2.3$ millones vinculados a los intereses futuros.

\solucion{13.3} El inventario es un activo reconocido cuyo \emph{valor} cambia con el precio del crudo: lo natural es cubrir su valor razonable y reconocer ambos cambios en resultados a la vez. El préstamo del BN genera \emph{flujos} futuros en soles cuyo equivalente en dólares es variable: se cubre la variabilidad de esos flujos, y el ORI permite diferir el efecto del derivado hasta que cada pago afecta el resultado. El NDF no necesita designarse porque las cuentas por pagar ya se ajustan por diferencia de cambio en resultados (NIC~21), igual que el NDF: la compensación ocurre sola en el mismo estado.

\solucion{13.4} El riesgo de sobrecobertura: si se cubriera más de lo que se tiene expuesto, la parte excedente sería una apuesta en la dirección contraria, es decir, especulación. Los límites inferiores evitan quedar descubierto; los superiores, especular. En la E2, el 20\,\% abierto protege además contra una baja del inventario.

\solucion{13.5} Umbral $=253.6-0.75\times60=253.6-45.0=208.6\to209$ millones; con tres bancos, unos 70 millones cada uno. Más líneas libres reducen el umbral que hay que pedir a los bancos.

\solucion{13.6} Porque así se evitan conflictos de interés y errores no detectados. Quien opera podría ocultar pérdidas o exceder límites si también los midiera o registrara. La separación garantiza que cada operación sea vista por alguien independiente, incluida la verificación de los precios que cotizan los bancos.

\solucion{13.7} Respuesta modelo: PETROPERÚ asume sin cobertura un pasivo de S/ 4,704 millones y un inventario de crudo de US\$ 268.6 millones, sin liquidez para absorber pérdidas. Recomendamos (1) un CCS sobre el préstamo del BN y NDF sobre el 80\,\% de las cuentas por pagar: reduce la pérdida ante una caída de 10\,\% del TC de US\$ 152.0 a 6.1 millones; cuesta el diferencial de tasas (0.24\,\% anual) y 11 pb de CVA; su riesgo residual es el cronograma inferido. (2) Un programa rodante sobre el 80\,\% del inventario, mitad swap y mitad collar de costo cero: reduce la pérdida ante una caída de 30\,\% del WTI de 80.6 a 34.0 millones sin prima; cede la subida y 3.7 millones de backwardation; su riesgo residual es la base Brent--WTI. Ambas se ejecutan solo con un umbral CSA de al menos US\$ 237 millones o la garantía del Estado, y con dispensa del incumplimiento cruzado por el CESCE.
